\subsection{Assignment 14 -- Datacube}
Per la creazione del cubo è stata configurata la
connessione al database tramite la definizione del Data Source.
Una volta stabilito il collegamento, abbiamo costruito il Data Source
View (DSV), nel quale abbiamo importato la fact table e tutte le
dimensioni rilevanti.
Per ciascuna dimensione è stato selezionato un sottoinsieme di
attributi, includendo esclusivamente quelli necessari per gli
assignment successivi.\\

All'interno della dimensione Date è stata definita una gerarchia
temporale, utilizzando gli attributi Year, Month e Day presenti nella
dimensione.
Per la definizione della gerarchia, sono stati introdotti due Named
Calculations:
\begin{itemize}
    \item \texttt{YearMonth} = concatenazione Year-Month nel formato YYYYMM
    \item \texttt{YearMonthDay} = concatenazione Year-Month-Day nel formato YYYYMMDD
\end{itemize}

Questi attributi ci hanno permesso di definire la gerarchia temporale:
\begin{center}
    Year → Month → Day
\end{center}
dove:
\begin{itemize}
    \item il livello Month utilizza come NameColumn l'attributo \texttt{YearMonth},
    \item il livello Day utilizza come NameColumn l'attributo \texttt{YearMonthDay}.
\end{itemize}

All'interno della dimensione Date era presente anche l'attributo Season,
mantenuto come attributo flat perché non introduce una relazione
di aggregazione naturale e stabile con gli altri livelli temporali
(giorno, mese, trimestre, anno).\\

All'interno del cubo abbiamo definito una seconda gerarchia, appartenente
alla dimensione Artist Geography, costruita per rappresentare il luogo
di nascita degli artisti:
\begin{center}
    Country → Region → Province → BirthPlace
\end{center}
Per ciascun livello di entrambe le gerarchie sono state mantenute le
versione sia l'attributo flat corrispondente.\\

Rispetto alla consegna del midterm, sono state apportate alcune
modifiche al modello:
\begin{itemize}
    \item eliminazione della misura \texttt{totalStreams},
    in quanto ridondante rispetto alla misura \texttt{stream1Month};
    \item eliminazione della misura calcolata \texttt{weightedStreams},
    il cui calcolo viene eseguito dalla rispettiva query.
\end{itemize}
In particolare, 

Inoltre, nel cubo è stata definita la dimensione \texttt{Role},
modellata come \emph{degenerate dimension} per l'attributo
\texttt{IsPrimary}. Tale attributo rappresenta una caratteristica
intrinseca del fatto, priva di ulteriori attributi descrittivi o
gerarchie.